\ifdefined\gkver\else\def\gkver{student}\fi
\documentclass[\gkver]{gaokaozhenti}
\begin{document}

\chapter{1996年全国}

\begin{enumerate}
\item
%% number: 1
%% paperName: 1996年全国高考第 1 题 4 分
%% typeId: 1
%% type: 单选题
%% chapter: 原子和原子核
%% point: 核聚变
%% method: 无
%% score: 4
%% degree: 850
%% duplicateId: 0
%% body:
下列核反应方程式中，表示核聚变过程的是\xzanswer{B}
\fourchoices[answer=B]
{\ce{^{30}_{15}P -> ^{30}_{14}Si + ^{0}_{1}e}}
{\ce{^{2}_{1}H + ^{3}_{1}H -> ^{4}_{2}He + ^{1}_{0}n}}
{\ce{^{14}_{6}C -> ^{14}_{7}N + ^{0}_{-1}e}}
{\ce{^{238}_{92}U -> ^{234}_{90}Th + ^{4}_{2}He}}
%% answer: B
%% memo:
\memoanswer{
核聚变是指多个较轻的核聚合成较重的核的反应，变化之前应至少有两种核，故 B 正确。
}

\item
%% number: 2
%% paperName: 1996年全国高考第 2 题 4 分
%% typeId: 1
%% type: 单选题
%% chapter: 光学
%% point: 光子说
%% method: 无
%% score: 4
%% degree: 850
%% duplicateId: 0
%% body:
红、橙、黄、绿四种单色光中，光子能量最小的是\xzanswer{A}
\fourchoices[answer=A]
{红光}
{橙光}
{黄光}
{绿光}
%% answer: A
%% memo:
\memoanswer{
由光子说知，光子的能量 $E = h\nu$，其能量跟频率成正比，而红光的频率最低，故 A 正确。
}

\item
%% number: 3
%% paperName: 1996年全国高考第 3 题 4 分
%% typeId: 1
%% type: 单选题
%% chapter: 电磁感应
%% point: 楞次定律
%% method: 无
%% score: 4
%% degree: 650
%% duplicateId: 0
%% body:
一平面线圈用细杆悬于 P 点，开始时细杆处于水平位置，释放后让它在如图所示的匀强磁场中运动。已知线圈平面始终与纸面垂直，当线圈第一次通过位置 Ⅰ 和位置 Ⅱ 时，顺着磁场的方向看去，线圈中感应电流的方向分别为\xzanswer{B}
\onepicture[w=5.3cm]{\includesvg{figs/03}}
\fourchoices[answer=B]
{逆时针方向　逆时针方向}
{逆时针方向　顺时针方向}
{顺时针方向　顺时针方向}
{顺时针方向　逆时针方向}
%% answer: B
%% memo:
\memoanswer{
当线圈第一次通过位置 Ⅰ 时，磁场方向向右通过线圈，且磁通量在增大，由楞次定律可确定感应电流方向为逆时针；同理，线圈第一次通过位置 Ⅱ 时，感应电流方向为顺时针，故 B 正确。
}

\item
%% number: 4
%% paperName: 1996年全国高考第 4 题 4 分
%% typeId: 1
%% type: 单选题
%% chapter: 气体性质
%% point: 阿伏伽德罗常数
%% method: 无
%% score: 4
%% degree: 850
%% duplicateId: 0
%% body:
只要知道下列哪一组物理量，就可以估算出气体中分子间的平均距离？\xzanswer{B}
\fourchoices[answer=B]
{阿伏伽德罗常数、该气体的摩尔质量和质量}
{阿伏伽德罗常数、该气体的摩尔质量和密度}
{阿伏伽德罗常数、该气体的质量和体积}
{该气体的密度、体积和摩尔质量}
%% answer: B
%% memo:
\memoanswer{
气体的摩尔体积 $=\frac{\text{摩尔质量}}{\text{密度}}$，一个气体分子所占有的空间等于摩尔体积除以阿伏加德罗常数，我们把每个分子所占有的空间近似看成正方体，每个正方体的边长即可看做分子之间的平均距离（认为正方体紧挨着排列），所以，已知摩尔质量与密度就可以算出摩尔体积，用摩尔体积除以阿伏加德罗常数就可算出每个分子所占有的空间，将其开三次方就可估算出分子的平均距离，故 B 正确。
}

\item
%% number: 5
%% paperName: 1996年全国高考第 5 题 4 分
%% typeId: 1
%% type: 单选题
%% chapter: 原子和原子核
%% point: 玻尔模型
%% method: 无
%% score: 4
%% degree: 850
%% duplicateId: 0
%% body:
根据玻尔理论，氢原子的电子由外层轨道跃迁到内层轨道后\xzanswer{D}
\fourchoices[answer=D]
{原子的能量增加，电子的动能减少}
{原子的能量增加，电子的动能增加}
{原子的能量减少，电子的动能减少}
{原子的能量减少，电子的动能增加}
%% answer: D
%% memo:
\memoanswer{
氢原子核外电子从离核较远的轨道跃迁到离核较近的轨道时，由 $E_{n} = \frac{1}{n^{2}}E_{1}$ 知，$E_{1}$ 为负值（$-13.6\UeV$），$n$ 减小，$E_{n}$ 减小；其动能可由 $k\frac{e^{2}}{r_{n}^{2}} = m\frac{v_{n}^{2}}{r_{n}}$ 得 $E_{kn} = \frac{1}{2}mv_{n}^{2} = \frac{ke^{2}}{2r_{n}}$，$r_{n}$ 减小，$E_{kn}$ 增大，故 D 正确。
}

\item
%% number: 6
%% paperName: 1996年全国高考第 6 题 4 分
%% typeId: 1
%% type: 单选题
%% chapter: 电路
%% point: 电容
%% method: 无
%% score: 4
%% degree: 650
%% duplicateId: 0
%% body:
在右图所示的实验装置中，平行板电容器的极板 A 与一灵敏的静电计相接，极板 B 接地。若极板 B 稍向上移动一点，由观察到的静电计指针变化作出平行板电容器电容变小的结论的依据是\xzanswer{D}
\onepicture[w=7.7cm]{figs/06.png}
\fourchoices[answer=D]
{两极板间的电压不变，极板上的电量变小}
{两极板间的电压不变，极板上的电量变大}
{极板上的电量几乎不变，两极板间的电压变小}
{极板上的电量几乎不变，两极板间的电压变大}
%% answer: D
%% memo:
\memoanswer{
静电计指针的变化表征了电容器两极板电势差的变化。由题意知，电容器带电量 $Q$ 不变，若极板 B 稍向上移动一点，可观察到电容器两极板电势差 $U$ 变大，由 $C = \frac{Q}{U}$ 可作出平行板电容器电容变小的结论，故 D 正确。
}

\item
%% number: 7
%% paperName: 1996年全国高考第 7 题 4 分
%% typeId: 1
%% type: 单选题
%% chapter: 光学
%% point: 透镜
%% method: 无
%% score: 4
%% degree: 650
%% duplicateId: 0
%% body:
一焦距为 $f$ 的凸透镜，主轴和水平的 $x$ 轴重合。$x$ 轴上有一光点位于透镜的左侧，光点到透镜的距离大于 $f$ 而小于 $2f$。若将此透镜沿 $x$ 轴向右平移 $2f$ 的距离，则在此过程中，光点经透镜所成的象点将\xzanswer{C}
\fourchoices[answer=C]
{一直向右移动}
{一直向左移动}
{先向左移动，接着向右移动}
{先向右移动，接着向左移动}
%% answer: C
%% memo:
\memoanswer{
凸透镜向右移动时，经历了在一倍焦距和二倍焦距之间、二倍焦距处、二倍焦距以外三个阶段。凸透镜成实像的三种情况：$u > 2f$，成倒立、缩小的实像，$f < v < 2f$；$u = 2f$，成倒立、等大的实像，$v = 2f$；$f < u < 2f$，成倒立、放大的实像，$v > 2f$。

初始时点光源到透镜的距离大于 $f$ 小于 $2f$，若将此透镜沿 $x$ 轴向右平移 $2f$ 的距离，点光源一直在凸透镜的一倍焦距以外，成倒立的实像。①凸透镜成实像时，点光源在二倍焦距处，物像之间的距离最小；②物体在凸透镜的一倍焦距和二倍焦距之间，在凸透镜的二倍焦距以外成倒立放大的实像；③凸透镜向右移动，凸透镜移动距离物体二倍焦距处，在凸透镜的二倍焦距处得到倒立等大的实像；④凸透镜继续向右移动，凸透镜距离物体在二倍焦距以外，在凸透镜的一倍焦距和二倍焦距之间成倒立缩小的实像；故成像先向左移动，然后再向右移动，故 C 正确。
}

\item
%% number: 8
%% paperName: 1996年全国高考第 8 题 4 分
%% typeId: 1
%% type: 单选题
%% chapter: 运动和力
%% point: 动量定理
%% method: 无
%% score: 4
%% degree: 650
%% duplicateId: 0
%% body:
质量为 $1.0\Ukg$ 的小球从高 $20\Um$ 处自由下落到软垫上，反弹后上升的最大高度为 $5.0\Um$。小球与软垫接触的时间为 $1.0\Us$，在接触时间内小球受到合力的冲量大小为\xzanswer{C}（空气阻力不计，$g$ 取 $10\Umsq$）。
\fourchoices[answer=C]
{$10\UNs$}
{$20\UNs$}
{$30\UNs$}
{$40\UNs$}
%% answer: C
%% memo:
\memoanswer{
小球与软垫作用的过程中，软垫对球的弹力是一个变力，设竖直向上的方向为正方向，由动量定理得 $I_{\text{合}} = mv_{2} - (-mv_{1}) = m(v_{2} + v_{1})$，对下落和弹起上升的过程，由机械能守恒得 $mgh = \frac{1}{2}mv^{2}$，刚与软垫接触时的速度 $v_{1} = \sqrt{2gh_{1}} = 20\Ums$，刚与软垫脱离时的速度 $v_{2} = \sqrt{2gh_{2}} = 10\Ums$，将 $v_{1}$、$v_{2}$ 的值分别代入上式得 $I_{\text{合}} = m(v_{2} + v_{1}) = 30\UNs$，故 C 正确。
}

\item
%% number: 9
%% paperName: 1996年全国高考第 9 题 6 分
%% typeId: 2
%% type: 多选题
%% chapter: 直线运动
%% point: 匀变速直线运动
%% method: 无
%% score: 6
%% degree: 550
%% duplicateId: 0
%% body:
一物体作匀变速直线运动，某时刻速度的大小为 $4\Ums$，1 秒钟后速度的大小变为 $10\Ums$。在这 1 秒钟内该物体的\xzanswer{AD}
\fourchoices[answer=AD]
{位移的大小可能小于 $4\Um$}
{位移的大小可能大于 $10\Um$}
{加速度的大小可能小于 $4\Umsq$}
{加速度的大小可能大于 $10\Umsq$}
%% answer: AD
%% memo:
\memoanswer{
物体做匀变速直线运动，有两种可能：①若物体做匀加速直线运动，则 $v_{t} = 10\Ums$，位移 $s = \frac{v_{0} + v_{t}}{2}t = 7\Um$，加速度 $a = \frac{v_{t} - v_{0}}{t} = 6\Umsq$；②若物体做匀减速直线运动，则 $v_{t} = -10\Ums$，位移 $s = \frac{v_{0} + v_{t}}{2}t = -3\Um$，加速度 $a = \frac{v_{t} - v_{0}}{t} = -14\Umsq$，故 AD 正确。
}

\item
%% number: 10
%% paperName: 1996年全国高考第 10 题 6 分
%% typeId: 2
%% type: 多选题
%% chapter: 交流电
%% point: LC振荡回路
%% method: 无
%% score: 6
%% degree: 650
%% duplicateId: 0
%% body:
LC 回路中电容两端的电压 $u$ 随时刻 $t$ 变化的关系如下图所示，则\xzanswer{BC}
\onepicture[w=5.8cm]{\includesvg{figs/10}}
\fourchoices[answer=BC]
{在时刻 $t_{1}$，电路中的电流最大}
{在时刻 $t_{2}$，电路的磁场能最大}
{从时刻 $t_{2}$ 至 $t_{3}$，电路的电场能不断增大}
{从时刻 $t_{3}$ 至 $t_{4}$，电容的带电量不断增大}
%% answer: BC
%% memo:
\memoanswer{
在时刻 $t_{1}$，电压最大，即电量最大，电路中电场能最大，电路中的电流为 0，故 A 错误；在时刻 $t_{2}$，电压为 0，即电量为 0，电路中的电流最大，电路的磁场能最大，B 正确；从时刻 $t_{2}$ 至 $t_{3}$，电压从 0 变到最大，即电量从 0 变到最大，电容器充电，电路的电场能不断增大，故 C 正确；从时刻 $t_{3}$ 至 $t_{4}$，电压从最大变到 0，即电量从最大变到 0，电容的带电量不断减小（电容器放电），故 D 错误。
}

\item
%% number: 11
%% paperName: 1996年全国高考第 11 题 6 分
%% typeId: 1
%% type: 单选题
%% chapter: 电场
%% point: 电场线
%% method: 无
%% score: 6
%% degree: 650
%% duplicateId: 0
%% body:
如图 a，b，c 是一条电场线上的三个点，电力线的方向由 a 到 c，a、b 间的距离等于 b、c 间的距离。用 $\varphi_{a}$、$\varphi_{b}$、$\varphi_{c}$ 和 $E_{a}$、$E_{b}$、$E_{c}$ 分别表示 a、b、c 三点的电势和电场强度，可以断定\xzanswer{A}
\onepicture[w=5.3cm]{\includesvg{figs/11}}
\fourchoices[answer=A]
{$\varphi_{a} > \varphi_{b} > \varphi_{c}$}
{$E_{a} > E_{b} > E_{c}$}
{$\varphi_{a} - \varphi_{b} = \varphi_{b} - \varphi_{c}$}
{$E_{a} = E_{b} = E_{c}$}
%% answer: A
%% memo:
\memoanswer{
不论是何种情况的电场，沿着电场线方向电势总是降低的，即 $\varphi_{a} > \varphi_{b} > \varphi_{c}$，故 A 正确；一根电场线不能判断电场线的疏密，所以不能判断电场强度的大小，故 BD 错误；由于不能判断电场强度的大小，也不能判断电势差的大小，故 C 错误。
}

\item
%% number: 12
%% paperName: 1996年全国高考第 12 题 6 分
%% typeId: 2
%% type: 多选题
%% chapter: 机械波
%% point: 波的多解性
%% method: 无
%% score: 6
%% degree: 450
%% duplicateId: 0
%% body:
一根张紧的水平弹性长绳上的 a、b 两点，相距 $14.0\Um$，b 点在 a 点的右方。当一列简谐横波沿此长绳向右传播时，若 a 点的位移达到正极大时，b 点的位移恰为零，且向下运动。经过 $1.00\Us$ 后，a 点的位移为零，且向下运动，而 b 点的位移恰达到负极大，则这简谐横波的波速可能等于\xzanswer{AC}
\onepicture[w=5.3cm]{\includesvg{figs/12}}
\fourchoices[answer=AC]
{$4.67\Ums$}
{$6\Ums$}
{$10\Ums$}
{$14\Ums$}
%% answer: AC
%% memo:
\memoanswer{
依题意画出波形图，如图所示。在 $t = 0\Us$ 时的波形图可能如图中实线所示，在 $t = 1.00\Us$ 时刻的波形图可能如图中虚线所示。由图可知 a、b 间至少相距 $\frac{3}{4}$ 个波长（图中画了 $\frac{3}{4}$ 个波长），由波的周期性知，a、b 间的距离 $ab = 14.0\Um = \frac{3}{4}\lambda + n_{1}\lambda$（$n_{1} = 0, 1, 2, \ldots$）；由图可知，$1.00\Us$ 内 a 点至少振动 $\frac{1}{4}$ 个周期，由波的周期性知，$t = 1.00\Us = \frac{1}{4}T + n_{2}T$（$n_{2} = 0, 1, 2, \ldots$），且 $n_{1}$、$n_{2}$ 是相互独立，解得 $\lambda = \frac{14.0}{\frac{3}{4} + n_{1}}$，$T = \frac{1.00}{\frac{1}{4} + n_{2}}$（$n_{1}$、$n_{2} = 0, 1, 2, \ldots$），则 $v = \frac{\lambda}{T} = \frac{14.0\times(1 + 4n_{2})}{3 + 4n_{1}}$。当 $n_{1} = 0$、$n_{2} = 0$ 时，$v_{0} = \frac{14}{3}\Ums \approx 4.67\Ums$；当 $n_{1} = 0$、$n_{2} = 1$ 时，$v_{1} = \frac{14\times5}{3}\Ums > 14\Ums$；当 $n_{1} = 1$、$n_{2} = 0$ 时，$v_{0} = \frac{14}{7}\Ums = 2\Ums$；当 $n_{1} = 1$、$n_{2} = 1$ 时，$v_{1} = \frac{14\times5}{7}\Ums = 10\Ums$；当 $n_{1} = 1$、$n_{2} = 2$ 时，$v_{2} = \frac{14\times9}{7}\Ums > 14\Ums$；当 $n_{1} = 2$、$n_{2} = 0$ 时，$v_{0} = \frac{14}{11}\Ums < 2\Ums$；当 $n_{1} = 2$、$n_{2} = 1$ 时，$v_{1} = \frac{14\times5}{11}\Ums \approx 6.36\Ums$；当 $n_{1} = 2$、$n_{2} = 2$ 时，$v_{2} = \frac{14\times9}{11}\Ums \approx 11.45\Ums$；当 $n_{1} = 2$、$n_{2} = 3$ 时，$v_{3} = \frac{14\times13}{11}\Ums > 14\Ums$。以此类推，得出无数个波速的可能值，只有 AC 符合，故 AC 正确。

\onepicture[w=7.9cm]{figs/12_an.png}
}

\item
%% number: 13
%% paperName: 1996年全国高考第 13 题 6 分
%% typeId: 2
%% type: 多选题
%% chapter: 运动和力
%% point: 动量守恒定律
%% method: 无
%% score: 6
%% degree: 650
%% duplicateId: 0
%% body:
半径相等的两个小球甲和乙，在光滑水平面上沿同一直线相向运动。若甲球的质量大于乙球的质量，碰撞前两球的动能相等，则碰撞后两球的运动状态可能是\xzanswer{AC}
\fourchoices[answer=AC]
{甲球的速度为零而乙球的速度不为零}
{乙球的速度为零而甲球的速度不为零}
{两球的速度均不为零}
{两球的速度方向均与原方向相反，两球的动能仍相等}
%% answer: AC
%% memo:
\memoanswer{
因两球半径相等且沿同一直线相向运动，可知碰后两球仍在这条直线上运动，即正碰。又因两球动能相等而质量不等，则 $m_{\text{甲}} > m_{\text{乙}}$，由 $p = \sqrt{2mE_{k}}$ 知，甲的动量 $p_{\text{甲}}$ 大于乙的动量 $p_{\text{乙}}$，系统的总动量不为零，总动量方向与甲的初动量方向相同。由动量守恒得 $p_{\text{甲}} - p_{\text{乙}} = p_{\text{甲}}' + p_{\text{乙}}'$。①

然后逐项研究：

A 项：若碰后甲球速度为 0，即 $v_{\text{甲}}' = 0$，由①式得 $p_{\text{乙}}' = p_{\text{甲}} - p_{\text{乙}} > 0$，$v_{\text{乙}}' > 0$，故 A 正确。

B 项：若碰后乙球速度为 0，即 $v_{\text{乙}}' = 0$，由①式得 $p_{\text{甲}}' = p_{\text{甲}} - p_{\text{乙}} > 0$，$v_{\text{甲}}' > 0$，甲球将超越乙沿原方向运动，这是不符合事实的，故 B 错误。

C 项：两球速度均不为 0，由①式可见，若两球碰撞后速度方向相同，$p_{\text{甲}}' > 0$，$p_{\text{乙}}' > 0$，①式成立且符合事实，故 C 正确。

D 项：若两球碰撞后速度反向，动能不变，则必有 $p_{\text{甲}}' = -p_{\text{甲}}$，$p_{\text{乙}}' = -p_{\text{乙}}$，代入①式中得 $p_{\text{甲}} - p_{\text{乙}} = -p_{\text{甲}} + p_{\text{乙}}$，即 $2p_{\text{甲}} = 2p_{\text{乙}}$，$p_{\text{甲}} = p_{\text{乙}}$，显然与已知条件 $p_{\text{甲}} > p_{\text{乙}}$ 矛盾，故 D 错误。
}

\item
%% number: 14
%% paperName: 1996年全国高考第 14 题 6 分
%% typeId: 2
%% type: 多选题
%% chapter: 机械振动
%% point: 简谐振动
%% method: 无
%% score: 6
%% degree: 450
%% duplicateId: 0
%% body:
如果表中给出的是作简谐振动的物体的位移 $x$ 或速度 $v$ 与时刻的对应关系，$T$ 是振动周期，则下列选项中正确的是\xzanswer{AB}

\begin{center}
\begin{tabular}{|c|c|c|c|c|c|}
\hline
 & $0$ & $T/4$ & $T/2$ & $3T/4$ & $T$ \\
\hline
甲 & 零 & 正向最大 & 零 & 负向最大 & 零 \\
\hline
乙 & 零 & 负向最大 & 零 & 正向最大 & 零 \\
\hline
丙 & 正向最大 & 零 & 负向最大 & 零 & 正向最大 \\
\hline
丁 & 负向最大 & 零 & 正向最大 & 零 & 负向最大 \\
\hline
\end{tabular}
\end{center}
\fourchoices[answer=AB]
{若甲表示位移 $x$，则丙表示相应的速度 $v$}
{若丁表示位移 $x$，则甲表示相应的速度 $v$}
{若丙表示位移 $x$，则甲表示相应的速度 $v$}
{若乙表示位移 $x$，则丙表示相应的速度 $v$}
%% answer: AB
%% memo:
\memoanswer{
如图所示为做简谐运动的物体的位移图象和速度图象（实线为位移图象，虚线为速度图象），由图可知，在平衡位置处，位移为零，速度为最大，在最大位移处，速度为 0，由图可知，CD 均不符合实际情况，故 AB 正确。

\onepicture[w=12cm]{figs/14_an.png}
}

\item
%% number: 15
%% paperName: 1996年全国高考第 15 题 5 分
%% typeId: 4
%% type: 实验题
%% chapter: 功和能
%% point: DIS研究机械能守恒定律
%% method: 无
%% score: 5
%% degree: 550
%% duplicateId: 0
%% body:
在“验证机械能守恒定律”的实验中，已知打点计时器所用电源的频率为 $50\UHz$。查得当地的重力加速度 $g = 9.80\Umsq$。测得所用的重物的质量为 $1.00\Ukg$。实验中得到一条点迹清晰的纸带，把第一个点记作 0，另选连续的 4 个点 A、B、C、D 作为测量的点。经测量知道 A、B、C、D 各点到 0 点的距离分别为 $62.99\Ucm$、$70.18\Ucm$、$77.76\Ucm$、$85.73\Ucm$。根据以上数据，可知重物由 0 点运动到 C 点，重力势能的减少量等于\tkanswer[1.8cm]{$7.62\UJ$}，动能的增加量等于\tkanswer[1.8cm]{$7.56\UJ$}（取 3 位有效数字）。
\onepicture[align=right, w=7.9cm]{figs/15.png}
%% answer: 
\jdanswer{
$7.62\UJ$，$7.56\UJ$
}
%% memo:
\memoanswer{
从 $O$ 点到 $C$ 点势能减少量 $\Delta E_{\text{p}} = mgh = 1.00\times9.80\times0.7776\UJ = 7.62\UJ$，动能的增加量 $\Delta E_{\text{k}} = \frac{1}{2}mv_{C}^{2}$，由匀变速直线运动一段时间内的平均速度等于这段时间中间时刻的瞬时速度，即有 $v_{C} = \frac{s_{BD}}{2T}$，又 $s_{BD} = s_{OD} - s_{OB} = (85.73 - 70.18)\Ucm = 15.55\Ucm$，$T = 0.02\Us$，故 $\Delta E_{\text{k}} = 7.56\UJ$。
}

\item
%% number: 16
%% paperName: 1996年全国高考第 16 题 6 分
%% typeId: 4
%% type: 实验题
%% chapter: 电场
%% point: 描绘等势面
%% method: 无
%% score: 6
%% degree: 650
%% duplicateId: 0
%% body:
在用电流场模拟静电场描绘电场等势线的实验中，在下列所给出的器材中，应该选用的是\tkanswer[1.4cm]{B、E}（用器材前的字母表示）。（A）$6\UV$ 的交流电源，（B）$6\UV$ 的直流电源，（C）$100\UV$ 的直流电源，（D）量程 $0\sim0.5\UV$、零刻度在刻度盘中央的电压表，（E）量程 $0\sim300\UuA$、零刻度在刻度盘中央的电流表。

\begin{enumerate}
	\item
	在实验过程中，要把复写纸、导电纸、白纸铺放在木板上，它们的顺序（自上而下）是①\tkanswer[1.6cm]{导电纸}②\tkanswer[1.6cm]{复写纸}③\tkanswer[1.6cm]{白纸}。
	\item
	在实验中，按下电键，接通电路。若一个探针与基准点 O 接触，另一探针已分别在基准点 O 的两侧找到了实验所需要的两点 a、b（如图所示），则当此探针与 a 点接触时，电表的指针应\tkanswer[1.4cm]{指零}（填“左偏”、“指零”或“右偏”）；当此探针与 b 点接触时，电表的指针应\tkanswer[1.2cm]{指零}（填“左偏”、“指零”或“右偏”）。
\end{enumerate}
\onepicture[align=right, w=6.6cm]{figs/16.png}
%% answer: 
\jdanswer{
\begin{enumerate}
	\item
	①导电纸②复写纸③白纸

	\item
	指零，指零

\end{enumerate}
}
%% memo:
\memoanswer{
要用直流电源形成稳恒电流场模拟静电场，应选 $6\UV$ 直流电源。要用探针寻找等势点，电流表灵敏度要高，指针能向左、向右偏转，导电纸应在最上面，涂有导电物质的面向上，以便用探针寻找等势点，白纸是为了记录等势点，应在下面，复写纸应在二者之间。
}

\item
%% number: 17
%% paperName: 1996年全国高考第 17 题 6 分
%% typeId: 4
%% type: 实验题
%% chapter: 电路
%% point: 伏安法测电阻
%% method: 无
%% score: 6
%% degree: 550
%% duplicateId: 0
%% body:
在用伏安法测电阻的实验中，所用电压表的内阻约为 $20\UkO$，电流表的内阻约为 $10\UO$，选择能够尽量减小误差的电路图接线进行实验，读得的各组数据用实心圆点标于坐标图上（如右图所示）。

\begin{enumerate}
	\item
	根据各点表示的数据描出 $I-U$ 图线，由此求得该电阻的阻值 $R_{x} =$\tkanswer[2.8cm]{$2.4\times10^{3}\UO$}（保留两位有效数字）。
	\item
	画出此实验的电路原理图。
\end{enumerate}
\onepicture[align=right, w=7.9cm]{figs/17.png}
%% answer: 
\jdanswer{
\begin{enumerate}
	\item
	舍去不合理点的直线，$R_{x} = 2.4\times10^{3}\UO$。

	\item
	如图所示。

	\onepicture[w=5.3cm]{figs/17_an2.png}

\end{enumerate}
}
%% memo:
\memoanswer{
$I = 0.8\UmA$ 的点偏离较大，故舍去，将其余各点连成直线，得到过坐标原点的 $U-I$ 图线，从图线上读取 $U$、$I$ 值，如取 $I = 1.0\UmA$，$U = 2.4\UV$，解得 $R_{x} = \frac{U}{I} = \frac{2.4}{1.0\times10^{-3}}\UO = 2.4\times10^{3}\UO$。由于待测电阻较大，$R_{x} \gg R_{\text{A}}$，测量时应取电流表内接法，滑动变阻器取分压接法。

\onepicture[w=7.9cm]{figs/17_an1.png}
}

\item
%% number: 18
%% paperName: 1996年全国高考第 18 题 5 分
%% typeId: 3
%% type: 填空题
%% chapter: 磁场
%% point: 磁力矩
%% method: 无
%% score: 5
%% degree: 650
%% duplicateId: 0
%% body:
如图所示，一细导体杆弯成四个拐角均为直角的平面折线，其 ab、cd 段长度均为 $l_{1}$，bc 段长度为 $l_{2}$。弯杆位于竖直平面内，Oa、dO$'$ 段由轴承支撑且沿水平放置。整个弯杆置于匀强磁场中，磁场方向竖直向上，磁感应强度为 $B$。今在导体杆中沿 abcd 通以大小为 $I$ 的电流，此时导体杆受到的安培力对 OO$'$ 轴的力矩大小等于\tkanswer[3cm]{$IBl_{1}l_{2}$}。
\onepicture[align=right, w=7.9cm]{figs/18.png}
%% answer: 
\jdanswer{
$IBl_{1}l_{2}$
}
%% memo:
\memoanswer{
$bc$ 边受安培力的大小为 $F = BIl_{2}$，以 $OO'$ 为轴，该安培力的力臂为 $l_{1}$，由力矩定义有 $M = BIl_{1}l_{2}$。
}

\item
%% number: 19
%% paperName: 1996年全国高考第 19 题 5 分
%% typeId: 7
%% type: 作图题
%% chapter: 电磁感应
%% point: 电磁感应中图象
%% method: 无
%% score: 5
%% degree: 550
%% duplicateId: 0
%% body:
右图中 abcd 为一边长为 $l$、具有质量的刚性导线框，位于水平面内，bc 边中串接有电阻 $R$，导线的电阻不计。虚线表示一匀强磁场区域的边界，它与线框的 ab 边平行。磁场区域的宽度为 $2l$，磁感应强度为 $B$，方向竖直向下。线框在一垂直于 ab 边的水平恒定拉力作用下，沿光滑水平面运动，直到通过磁场区域。已知 ab 边刚进入磁场时，线框便变为匀速运动，此时通过电阻 $R$ 的电流的大小为 $i_{0}$，试在右图的 $i-x$ 坐标上定性画出：从导线框刚进入磁场到完全离开磁场的过程中，流过电阻 $R$ 的电流 $i$ 的大小随 ab 边的位置坐标 $x$ 变化的曲线。
\onepicture[align=right, w=7.9cm]{figs/19a.png}
\drawpicanswer{\onepicture[align=right, w=9cm]{figs/19b.png}}{\onepicture[align=right, w=10.6cm]{figs/19_an.png}}
%% answer: 
\jdanswer{
电流 $i$ 随 $x$ 的变化规律如下表所示，变化曲线如图所示。

\begin{center}
\begin{tabular}{|c|c|}
\hline
$x$ & $i$ \\
\hline
$0\to l$ & $i = i_{0}$ \\
\hline
$l\to 2l$ & $i = 0$ \\
\hline
$2l\to 3l$ & 如下图，且 $x = 3l$ 处 $i \geqslant i_{0}$ \\
\hline
\end{tabular}
\end{center}
}
%% memo:
\memoanswer{
由题意，ab 边刚进入磁场，线框便匀速运动，所以在线框匀速运动 $l$ 距离过程中回路电流恒为 $i_{0}$；当线框在 $x = l$ 至 $x = 2l$ 间运动时回路无磁通变化，故 $i = 0$，但该段距离内线框受到合外力为水平拉力，所以线框做匀加速运动；当 ab 边出磁场时（$x = 2l$），线框的速度大于刚进磁场时的速度，cd 段切割磁感线产生的电流 $i > i_{0}$，同时受到的安培力大于 $F$，做减速运动，随着速度的减小，安培力变小，减速运动的加速度变小，电流的变化变慢，但 cd 边出磁场时（$x = 3l$）最小速度不能小于进磁场时的速度，所以 $x = 3l$ 处回路中电流不能小于 $i_{0}$。
}

\item
%% number: 20
%% paperName: 1996年全国高考第 20 题 5 分
%% typeId: 3
%% type: 填空题
%% chapter: 力物体的平衡
%% point: 物体系平衡
%% method: 无
%% score: 5
%% degree: 450
%% duplicateId: 0
%% body:
如图所示，倔强系数为 $k_{1}$ 的轻质弹簧两端分别与质量为 $m_{1}$、$m_{2}$ 的物块 1、2 拴接，倔强系数为 $k_{2}$ 的轻质弹簧上端与物块 2 拴接，下端压在桌面上（不拴接），整个系统处于平衡状态。现施力将物块 1 缓慢竖直上提，直到下面那个弹簧的下端刚脱离桌面。在此过程中，物块 2 的重力势能增加了\tkanswer[2.8cm]{$\frac{m_{2}(m_{1} + m_{2})g^{2}}{k_{2}}$}，物块 1 的重力势能增加了\tkanswer[4cm]{$m_{1}(m_{1} + m_{2})g^{2}\left(\frac{1}{k_{1}} + \frac{1}{k_{2}}\right)$}。
\onepicture[align=right, w=2.1cm]{figs/20.png}
%% answer: 
\jdanswer{
$\frac{m_{2}(m_{1} + m_{2})g^{2}}{k_{2}}$，$m_{1}(m_{1} + m_{2})g^{2}\left(\frac{1}{k_{1}} + \frac{1}{k_{2}}\right)$
}
%% memo:
\memoanswer{
设上面弹簧为弹簧 1，下面弹簧为弹簧 2，在原题图所示状态，将物体 1 和物体 2 视为一个整体，弹簧 2 被压缩 $\Delta x_{2}$，由胡克定律得 $F_{2} = k_{2}\Delta x_{2} = (m_{1} + m_{2})g$，解得 $\Delta x_{2} = \frac{(m_{1} + m_{2})g}{k_{2}}$，即原题图所示状态，弹簧 2 比原长缩短了 $\Delta x_{2}$；当把物块 1 缓慢上提到弹簧 2 的下端刚脱离桌面时，弹簧 2 应为原长，由此可知，物块 2 上升的高度为 $\Delta x_{2}$，所以物块 2 的重力势能增加了 $\Delta E_{\text{p}2} = m_{2}g\Delta x_{2} = \frac{m_{2}(m_{1} + m_{2})g^{2}}{k_{2}}$。

在原题图所示状态，弹簧 1 被压缩 $\Delta x_{1}$，弹簧 1 的弹力 $F_{1} = k_{1}\Delta x_{1} = m_{1}g$，解得 $\Delta x_{1} = \frac{m_{1}g}{k_{1}}$；当把物块 1 缓慢上提到弹簧 2 的下端刚脱离桌面时，因为物块 2 是因弹簧 1 的拉力与重力保持平衡，所以弹簧 1 被拉伸 $\Delta x_{1}'$，弹簧 1 的弹力 $F_{1}' = k_{1}\Delta x_{1}' = m_{2}g$，解得 $\Delta x_{1}' = \frac{m_{2}g}{k_{1}}$，即弹簧 1 由被压缩 $\Delta x_{1}$ 到被拉伸 $\Delta x_{1}'$，弹簧 1 的伸长量变化了 $\Delta x_{\text{总}1} = \frac{(m_{1} + m_{2})g}{k_{1}}$，即物块 1 上升了 $\frac{(m_{1} + m_{2})g}{k_{2}} + \frac{(m_{1} + m_{2})g}{k_{1}}$，所以物块 1 的重力势能增加了 $\Delta E_{\text{p}1} = m_{1}g\Delta h = m_{1}g\left(\frac{m_{1} + m_{2}}{k_{1}}g + \frac{m_{1} + m_{2}}{k_{2}}g\right) = m_{1}(m_{1} + m_{2})g^{2}\left(\frac{1}{k_{1}} + \frac{1}{k_{2}}\right)$。
}

\item
%% number: 21
%% paperName: 1996年全国高考第 21 题 5 分
%% typeId: 3
%% type: 填空题
%% chapter: 功和能
%% point: 动能定理
%% method: 无
%% score: 5
%% degree: 350
%% duplicateId: 0
%% body:
在光滑水平面上有一静止的物体。现以水平恒力甲推这一物体，作用一段时间后，换成相反方向的水平恒力乙推这一物体。当恒力乙作用时间与恒力甲作用时间相同时，物体恰好回到原处，此时物体的动能为 $32\UJ$，则在整个过程中，恒力甲做的功等于\tkanswer[1.4cm]{$8\UJ$}，恒力乙做的功等于\tkanswer[1.4cm]{$24\UJ$}。
%% answer: 
\jdanswer{
$8\UJ$，$24\UJ$
}
%% memo:
\memoanswer{
设水平恒力甲为 $F_{1}$，作用时间 $t$，水平恒力乙为 $F_{2}$，物体质量为 $m$，由甲、乙两力作用下物体运动的位移关系得 $-\frac{1}{2}\cdot\frac{F_{1}}{m}t^{2} = \frac{F_{1}}{m}\cdot t^{2} - \frac{1}{2}\cdot\frac{F_{2}}{m}t^{2}$，解得 $F_{2} = 3F_{1}$，恒力甲推物体做的功为 $W_{1} = F_{1}s$，恒力乙推物体做的功为 $W_{2} = F_{2}s$，联立解得 $W_{2} = 3W_{1}$ ①，由动能定理得 $W_{1} + W_{2} = E_{k2} - 0 = 32\UJ$ ②，由①②解得 $W_{1} = 8\UJ$，$W_{2} = 24\UJ$。
}

\item
%% number: 22
%% paperName: 1996年全国高考第 22 题 5 分
%% typeId: 6
%% type: 计算题
%% chapter: 运动和力
%% point: 牛顿定律的应用
%% method: 无
%% score: 5
%% degree: 850
%% duplicateId: 0
%% body:
一物块从倾角为 $\theta$、长为 $s$ 的斜面的顶端由静止开始下滑，物块与斜面的滑动摩擦系数为 $\mu$，求物块滑到斜面底端所需的时间。
%% answer: 
\jdanswer{
$t = \sqrt{\frac{2s}{g\sin\theta - \mu g\cos\theta}}$
}
%% memo:
\memoanswer{
设物块质量为 $m$，加速度为 $a$，物块受力情况如图所示，$mg\sin\theta - f = ma$，$N - mg\cos\theta = 0$，$f = \mu N$，解得 $a = g\sin\theta - \mu g\cos\theta$，由 $s = \frac{1}{2}at^{2}$ 得 $t = \sqrt{\frac{2s}{a}} = \sqrt{\frac{2s}{g\sin\theta - \mu g\cos\theta}}$。

\onepicture[w=4.2cm]{figs/22_an.png}
}

\item
%% number: 23
%% paperName: 1996年全国高考第 23 题 8 分
%% typeId: 6
%% type: 计算题
%% chapter: 光学
%% point: 光的折射
%% method: 无
%% score: 8
%% degree: 550
%% duplicateId: 0
%% body:
在折射率为 $n$、厚度为 $d$ 的玻璃平板上方的空气中有一点光源 S，从 S 发出的光线 SA 以角度 $\theta$ 入射到玻璃板上表面，经过玻璃板后从下表面射出，如右图所示。若沿此光线传播的光从光源到玻璃板上表面的传播时间与在玻璃板中的传播时间相等，点光源 S 到玻璃上表面的垂直距离 $l$ 应是多少？
\onepicture[align=right, w=5.3cm]{figs/23.png}
%% answer: 
\jdanswer{
$l = \frac{n\cos\theta}{\sqrt{1 - \frac{1}{n^{2}}\sin^{2}\theta}}d$
}
%% memo:
\memoanswer{
设光线在玻璃中的折射角为 $r$，则光线从 S 到玻璃板上表面的传播距离为 $s_{1} = \frac{l}{\cos\theta}$ ①，光线从 S 到玻璃板上表面的传播时间为 $t_{1} = \frac{l}{c\cos\theta}$ ②，其中 $c$ 表示空气中的光速。光线在玻璃板中的传播距离为 $s_{2} = \frac{d}{\cos r}$ ③，光线在玻璃板中的传播时间为 $t_{2} = \frac{nd}{c\cos r}$ ④，由题意得 $t_{1} = t_{2}$，即 $\frac{nd}{c\cos r} = \frac{l}{c\cos\theta}$ ⑤，由折射定律 $\sin\theta = n\sin r$ ⑥，解得 $l = \frac{n\cos\theta}{\sqrt{1 - \frac{1}{n^{2}}\sin^{2}\theta}}d$ ⑦。
}

\item
%% number: 23
%% paperName: 1996年全国高考第 23 题 8 分
%% typeId: 6
%% type: 计算题
%% chapter: 运动和力
%% point: 动量和能量
%% method: 无
%% score: 8
%% degree: 450
%% duplicateId: 0
%% body:
一质量为 $M$ 的长木板静止在光滑水平桌面上。一质量为 $m$ 的小滑块以水平速度 $v_{0}$ 从长木板的一端开始在木板上滑动，直到离开木板。滑块刚离开木板时的速度为 $\frac{v_{0}}{3}$。若把该木板固定在水平桌面上，其它条件相同，求滑块离开木板时的速度 $v$。
%% answer: 
\jdanswer{
$v = \frac{v_{0}}{3}\sqrt{1 + \frac{4m}{M}}$
}
%% memo:
\memoanswer{
设第一次滑块离开时木板速度为 $v_{\text{木}}$，由系统的动量守恒，有 $mv_{0} = m\frac{v_{0}}{3} + Mv_{\text{木}}$ ①，设滑块与木板间摩擦力为 $f$，木板长 $L$，滑行距离 $s$，如图，由动能定理，对木板：$fs = \frac{1}{2}Mv_{\text{木}}^{2}$ ②，对滑块：$f(L + s) = \frac{1}{2}mv_{0}^{2} - \frac{1}{2}m\left(\frac{v_{0}}{3}\right)^{2}$ ③，当板固定时，有 $fL = \frac{1}{2}mv_{0}^{2} - \frac{1}{2}mv^{2}$ ④，解得 $v = \frac{v_{0}}{3}\sqrt{1 + \frac{4m}{M}}$。

\onepicture[w=7.9cm]{figs/23_an.png}
}

\item
%% number: 24
%% paperName: 1996年全国高考第 24 题 12 分
%% typeId: 6
%% type: 计算题
%% chapter: 气体性质
%% point: 玻意耳定律
%% method: 分情况讨论
%% score: 12
%% degree: 350
%% duplicateId: 0
%% body:
如图所示，有一个直立的气缸，气缸底到气缸口的距离为 $L_{0}\Ucm$，用一厚度和质量均可忽略不计的刚性活塞 A，把一定质量的空气封在气缸内，活塞与气缸间的摩擦可忽略。平衡时活塞上表面与气缸口的距离很小（计算时可忽略不计），周围大气的压强为 $H_{0}\UcmHg$。现把盛有水银的一个瓶子放在活塞上（瓶子的质量可忽略），平衡时活塞到气缸底的距离为 $L\Ucm$。若不是把这瓶水银放在活塞上，而是把瓶内水银缓缓不断地倒在活塞上方，这时活塞向下移，压缩气体，直到活塞不再下移。求此时活塞在气缸内可能的位置以及与之相对应的条件（即题中给出量之间应满足的关系）。设气体的温度不变。
\onepicture[align=right, w=3.7cm]{figs/24.png}
%% answer: 
\jdanswer{
\begin{enumerate}
	\item
	若 $L \geqslant H_{0}$，则 $L' = L$。

	\item
	若 $L < H_{0}$，则 $L' = H_{0}$。

\end{enumerate}
}
%% memo:
\memoanswer{
设整瓶水银放在活塞上后，使气缸内气体增加的压强为 $h\UcmHg$，由压强平衡得 $p_{1} = H_{0} + h$，由气体的初末状态知，气体做等温变化，由玻意耳定律得 $H_{0}L_{0} = (H_{0} + h)L$ ①，解得 $h = \frac{H_{0}(L_{0} - L)}{L}$ ②，$h$ 的大小反映了水银质量的大小。

当水银注入后，活塞不再下移时，设活塞上水银的深度为 $\Delta H\Ucm$，活塞下移的距离为 $\Delta x\Ucm$，气缸内气体的压强为 $p_{2}$，由压强平衡得 $p_{2} = H_{0} + \Delta H$，由玻意耳定律得 $H_{0}L_{0} = (H_{0} + \Delta H)(L_{0} - \Delta x)$ ③，解得 $\Delta H = \frac{H_{0}}{L_{0} - \Delta x}\Delta x$ ④。

可能发生两种情况：

1. 水银比较少，瓶内水银全部注入后，尚未灌满或刚好灌满活塞上方的气缸，这时气体增加的压强也为 $h\UcmHg$，则有 $\Delta H = h$ ⑤，$\Delta H \leqslant \Delta x$ ⑥，由（2）、（4）、（5）三式，得 $\Delta x = L_{0} - L$ ⑦，活塞到气缸底的距离 $L' = L_{0} - \Delta x = L$ ⑧，由（4）、（6）、（7）三式，得 $L \geqslant H_{0}$ ⑨，即若 $L \geqslant H_{0}$，则 $L' = L$。

2. 瓶内水银比较多，当活塞上方的气缸灌满水银时，瓶内还剩有一定量的水银，这时 $\Delta H = \Delta x$（10），$\Delta H < h$（11），由（4）、（10）两式，得 $\Delta x = L_{0} - H_{0}$（12），活塞到气缸底的距离 $L' = L_{0} - \Delta x = H_{0}$（13），由（2）、（10）、（11）三式，得 $L < H_{0}$（14），即若 $L < H_{0}$，则 $L' = H_{0}$。
}

\item
%% number: 25
%% paperName: 1996年全国高考第 25 题 12 分
%% typeId: 6
%% type: 计算题
%% chapter: 磁场
%% point: 带电粒子在复合场中的运动
%% method: 无
%% score: 12
%% degree: 250
%% duplicateId: 0
%% body:
设在地面上方的真空室内存在匀强电场和匀强磁场。已知电场强度和磁感应强度的方向是相同的，电场强度的大小 $E = 4.0\UVm$，磁感应强度的大小 $B = 0.15\UT$。今有一个带负电的质点以 $v = 20\Ums$ 的速度在此区域内沿垂直场强方向做匀速直线运动，求此带电质点的电量与质量之比 $\frac{q}{m}$ 以及磁场的所有可能方向（角度可用反三角函数表示）。
%% answer: 
\jdanswer{
\begin{enumerate}
	\item
	$\frac{q}{m} = 1.96\UCkg$。

	\item
	磁场是沿着与重力方向夹角 $\theta = \arctan 0.75$，且斜向下方的一切方向。

\end{enumerate}
}
%% memo:
\memoanswer{
根据带电质点做匀速直线运动的条件，得知此带电质点所受的重力、电场力和洛伦兹力的合力必定为零。由此推知此三个力在同一竖直平面内，如右图所示，质点的速度垂直纸面向外。

解法一：由合力为零的条件，可得 $mg = q\sqrt{(vB)^{2} + E^{2}}$ ①，求得带电质点的电量与质量之比 $\frac{q}{m} = \frac{g}{\sqrt{(vB)^{2} + E^{2}}}$ ②，代入数据得 $\frac{q}{m} = \frac{9.8}{\sqrt{(20\times0.15)^{2} + 4.0^{2}}}\UCkg = 1.96\UCkg$ ③，因质点带负电，电场方向与电场力方向相反，因而磁场方向也与电场力方向相反。设磁场方向与重力方向之间夹角为 $\theta$，则有 $qE\sin\theta = qvB\cos\theta$，解得 $\tan\theta = \frac{vB}{E} = \frac{20\times0.15}{4.0} = 0.75$，$\theta = \arctan 0.75$ ④，即磁场是沿着与重力方向夹角 $\theta = \arctan 0.75$，指向斜向下方。

解法二：因质点带负电，电场方向与电场力方向相反，因而磁场方向也与电场力方向相反。设磁场方向与重力方向间夹角为 $\theta$，由合力为零的条件，可得 $qE\sin\theta = qvB\cos\theta$ ①，$qE\cos\theta + qvB\sin\theta = mg$ ②，解得 $\frac{q}{m} = \frac{g}{\sqrt{(vB)^{2} + E^{2}}}$ ③，代入数据得 $\frac{q}{m} = 1.966\UCkg$ ④，$\tan\theta = \frac{vB}{E} = 0.75$，$\theta = \arctan 0.75$ ⑤，即磁场是沿着与重力方向成夹角 $\theta = \arctan 0.75$，指向斜向下方。

\onepicture[w=5.3cm]{figs/25_an.png}
}

\end{enumerate}

\end{document}
